\documentclass{article}
\usepackage{lmodern}
\usepackage{amsmath}
\usepackage{listings}
\usepackage{fullpage}
\usepackage{parskip}
\usepackage{xcolor}
\lstset{
	basicstyle=\ttfamily,
	columns=fullflexible,
	frame=single,
	breaklines=true,
	postbreak=\mbox{\textcolor{red}{$\hookrightarrow$}\space},
}
\begin{document}
\title{Experiment `p\_6\_2\_d\_replace\_with\_larger\_neighbor\_array' Results}
\author{\today}
\date{}
\maketitle
\textbf{Experiment outcome: }SUCCESS
\\\textbf{Bad responses: }0
\\\textbf{Responses containing}~\texttt{assume}~\textbf{: }0
\\\textbf{Responses containing}~\texttt{expect}~\textbf{: }0
\\\textbf{Resolution attempts: }1
\\\textbf{Hard fails (resolution): }0
\\\textbf{Soft fails (resolution): }0
\\\textbf{Verification attempts: }1
\section*{Problem Specification}
\textbf{Problem name: }p\_6\_2\_d\_replace\_with\_larger\_neighbor\_array
\\\textbf{Natural language statement: }Write a method to replace each element of an array, except the first and last, by the larger of its two neighbors.
\\\textbf{Method signature: }p\_6\_2\_d\_replace\_with\_larger\_neighbor\_array(inputs: array<int>) returns (result: array<int>)
\subsection*{Ensures}
\begin{itemize}
\item \texttt{result.Length == inputs.Length}
\item \texttt{result[0] == inputs[0]}
\item \texttt{result[inputs.Length - 1] == inputs[inputs.Length - 1]}
\item \texttt{forall i :: 1 <= i < inputs.Length - 1 ==> result[i] == if inputs[i - 1] >= inputs[i + 1] then inputs[i - 1] else inputs[i + 1]}
\end{itemize}
\subsection*{Requires}
\begin{itemize}
\item \texttt{inputs.Length >= 3}
\end{itemize}
\clearpage
\section*{GenAI interactions}
Below you will find all interactions between the `user' (program) and the `assistant' (OpenAI).
\subsection*{Program $\rightarrow$ GenAI}
\begin{lstlisting}
You are given the following task to perform in Dafny:

Write a method to replace each element of an array, except the first and last, by the larger of its two neighbors.

The signature should be:

p_6_2_d_replace_with_larger_neighbor_array(inputs: array<int>) returns (result: array<int>)

The method should respect the following contract:

requires inputs.Length >= 3, ensures result.Length == inputs.Length, ensures result[0] == inputs[0], ensures result[inputs.Length - 1] == inputs[inputs.Length - 1], ensures forall i :: 1 <= i < inputs.Length - 1 ==> result[i] == if inputs[i - 1] >= inputs[i + 1] then inputs[i - 1] else inputs[i + 1]

Produce and show only the Dafny body of this method, including the curly braces that surround it. Do not show the signature nor contract. You must not use 'assume' nor 'expect' anywhere in your code.
\end{lstlisting}
\subsection*{GenAI $\rightarrow$ Program}
\textbf{System fingerprint: }null
\\\textbf{ID: }chatcmpl-Dn0GhbJ3ywaU7WLf8g7cf1JzcpWVM
\\\textbf{Model: }gpt-5.5-2026-04-23
\\\textbf{Created at: }1780571475
\\\textbf{Finish reason: }stop
\\\textbf{Usage: }454 tokens in, and 747 tokens out
\begin{lstlisting}
{
  result := new int[inputs.Length];
  var n := inputs.Length;

  result[0] := inputs[0];
  result[n - 1] := inputs[n - 1];

  var i := 1;
  while i < n - 1
    invariant n == inputs.Length
    invariant result.Length == n
    invariant 1 <= i <= n - 1
    invariant result[0] == inputs[0]
    invariant result[n - 1] == inputs[n - 1]
    invariant forall j :: 1 <= j < i ==> result[j] == if inputs[j - 1] >= inputs[j + 1] then inputs[j - 1] else inputs[j + 1]
  {
    result[i] := if inputs[i - 1] >= inputs[i + 1] then inputs[i - 1] else inputs[i + 1];
    i := i + 1;
  }
}
\end{lstlisting}
\section*{Final Program}
\begin{lstlisting}
method p_6_2_d_replace_with_larger_neighbor_array(inputs: array<int>) returns (result: array<int>)
	requires inputs.Length >= 3
	ensures result.Length == inputs.Length
	ensures result[0] == inputs[0]
	ensures result[inputs.Length - 1] == inputs[inputs.Length - 1]
	ensures forall i :: 1 <= i < inputs.Length - 1 ==> result[i] == if inputs[i - 1] >= inputs[i + 1] then inputs[i - 1] else inputs[i + 1]
{
  result := new int[inputs.Length];
  var n := inputs.Length;

  result[0] := inputs[0];
  result[n - 1] := inputs[n - 1];

  var i := 1;
  while i < n - 1
    invariant n == inputs.Length
    invariant result.Length == n
    invariant 1 <= i <= n - 1
    invariant result[0] == inputs[0]
    invariant result[n - 1] == inputs[n - 1]
    invariant forall j :: 1 <= j < i ==> result[j] == if inputs[j - 1] >= inputs[j + 1] then inputs[j - 1] else inputs[j + 1]
  {
    result[i] := if inputs[i - 1] >= inputs[i + 1] then inputs[i - 1] else inputs[i + 1];
    i := i + 1;
  }
}
\end{lstlisting}
\section*{Total Token Usage}
\textbf{Input tokens: }454
\\\textbf{Output tokens: }747
\\\textbf{Reasoning tokens: }512
\\\textbf{Sum of `total tokens': }1201
\section*{Experiment Timings}
\textbf{Overall Experiment} started at 1780571474551, ended at 1780571518673, lasting 44122ms (44.12 seconds)
\\\textbf{Iteration \#1} started at 1780571474553, ended at 1780571518673, lasting 44120ms (44.12 seconds)
\end{document}
